\begin{solapartat}
Si la intensitat de camp gravitatori és uniforme, el mòdul de la diferència de potencial s'expressa com:

\punts{0,3} $|\Delta V| = gd$, llavors:

\punts{0,25} $g = \dfrac{|\Delta V|}{d} = \dfrac{74,4\ \mathrm{J/kg}}{20,0\ \mathrm{m}} = 3,72\ \mathrm{m/s^2}$

No es pot considerar com a correcta cap deducció que arribi a la relació $g = G\dfrac{M}{R^2}$ atès que es demana explícitament que es determini $g$ a partir de la diferència de potencial. L'objectiu d'aquest problema és comprovar que l'estudiant coneix i sap aplicar la relació existent entre el camp i la diferència de potencial quan el camp és uniforme.

\punts{0,4} El potencial gravitatori és una magnitud escalar, per tant, no poden ser ni l'esquema ``a'' ni l'esquema ``b''. D'altra banda, les línies de camp gravitatori indiquen la direcció de la força que actua sobre una massa puntual, per tant, són perpendiculars a la superfície del planeta i van dirigides de dalt cap a baix (línies vermelles a l'esquema). Com que les superfícies equipotencials són perpendiculars a les línies de camp, aquestes són paral·leles a la superfície, i, per tant, en la projecció al pla del dibuix corresponen a les línies horitzontals, esquema ``c''.

\punts{0,3} Les línies de camp indiquen la direcció en la qual el potencial disminueix, per tant, el potencial serà més gran com més allunyats estiguem de la superfície. Una justificació alternativa és a partir de l'expressió del potencial $V(r) = -G\dfrac{M}{r}$, atès que quan ens allunyem de la superfície $r$ disminueix, llavors el potencial augmenta.

\figura[0.35]{sol_fig1.png}

Tot aquest apartat també es pot resoldre partint de l'expressió (vàlida quan $g$ és uniforme):
\[ U = mgh \Rightarrow V = \frac{U}{m} = gh, \qquad \text{on } U_0 = U(h=0) = 0\ \mathrm{J} \]
Aquesta darrera expressió ens permet determinar $g$, i justificar com són les línies de camp i que $V$ augmenta quan $h$ augmenta.
\end{solapartat}

\begin{solapartat}
\punts{0,5} El treball que fa el camp gravitatori és:
\[ W_{cg} = -\Delta U = -m\Delta V = -1025(-74,4) = 76300\ \mathrm{J} \]

Observeu que $\Delta V = -74,4\ \mathrm{J/kg}$ atès que partim d'un potencial alt per acabar en un potencial més baix, per tant, la diferència de potencial és negativa. Alternativament:
\[ W_{cg} = -\Delta U = -(U_{fin} - U_{in}) = -(0 - mgh) = mgh = 76300\ \mathrm{J} \]

\punts{0,5} D'altra banda el teorema de les forces vives estableix que el treball total és igual a la variació d'energia cinètica. Com que la sonda baixa a velocitat constant, la variació d'energia cinètica és zero i també ho serà el treball total.

\punts{0,25} Com que el treball fet per les forces de fregament és nul, el treball total és la suma del treball fet pel camp gravitatori i el treball fet per la grua,
\[ W_{grua} + W_{cg} = 0 \Rightarrow W_{grua} = -W_{cg} = -76300\ \mathrm{J} \]

Per obtenir aquests 0,25 p, no només cal donar el valor correcte del treball, cal que el signe sigui també el correcte. La força que fa la grua s'oposa al moviment, per tant, aquest treball és negatiu.

Alternativament, tot aquest apartat es pot resoldre calculant el treball fet per la grua com la força pel desplaçament, tenint en compte que la força que fa la grua és $-mg$.
\end{solapartat}
